\section{Christian Gnosis}
\label{sec:ChristianGnosis}

In which we review \emph{Christian Gnosis: From St. Paul to Meister Eckhart} by \textbf{Wolfgang Smith}.

\subsection{Introduction}

\begin{quotex}
\textbf{St. Thomas Aquinas} was not the only one. Just as he arrived at contemplation through scholastic reasoning, so did the peak of the scholastic wave reach gnosis [\emph{mystique}], that is to say, intuition or the state of union of faith and intelligence, which is the aim of scholasticism. A \textbf{Meister Eckhart}, a \textbf{Ruysbroeck}, the Admirable Doctor, a \textbf{St. John of the Cross} are in fact spirits amongst whom you will search in vain for a spirit of opposition to scholasticism. For them also it was true that scholasticism was ``like straw", but they knew at the same time from their own experience that this straw is an excellent combustible. They certainly surpassed scholasticism, but after having attained its aim. For the aim of scholastic effort is contemplation, and it is gnosis [\emph{mystique}] which is the fruit of the scholastic tree. \flright{\textsc{Valentin Tomberg}, \emph{Letter XIX The Sun}}

\end{quotex}
The idea of a Christian gnosis is not widely accepted. That is as it should be, since gnosis is not something that can be taught but must be experienced. It is not necessary for salvation, so it may be safely ignored if it means nothing to you, and it should be ignored. Nevertheless, there are some who are called to follow the path of gnosis.

If you are not convinced of the existence of a Christian gnosis, there are resources you can turn to. \emph{Knowledge \& the Sacred} by \textbf{Seyyed Hossein Nasr} is a near encyclopedic compendium of gnosis. \emph{The Transcendental Universe} by \textbf{C. G. Harrison} deals with specific issues of Christian Gnosis. Finally, \emph{Christianity: Dogmatic Faith \& Gnostic Vivifying Knowledge} by \textbf{George Heart} deals with the transmutation of faith to gnosis.

Here, we will review \emph{Christian Gnosis: From St. Paul to Meister Eckhart} by \textbf{Wolfgang Smith}. Unlike other men of Tradition who accept a more or less universal gnosis, Smith will show that a specifically Christian gnosis goes beyond what other sapiential traditions can offer. He bases this on the doctrine of the Trinity, especially as it appears in Meister Eckhart.

We can but point out the highlights of Dr. Smith's book, and we do it with the following goals in mind:

\begin{itemize}
\item To encourage you to read Dr. Smith's book yourself 
\item To encourage you to read the saints and sages mentioned by Smith 
\item To encourage you to follow the path of gnosis yourself, if you are called to it 
\end{itemize}
\paragraph{Paul}
\begin{quotex}
Woe to you lawyers, for you have taken away the key of gnosis: you yourselves have not entered in, and those that were entering in, you have hindered. \flright{\textsc{Luke 11:52}}

\end{quotex}
Smith announces his intent

\begin{quotex}
is to understand, as clearly as we can, what gnosis is, and what role it plays in the perfection of Christian life.

\end{quotex}
First of all, he justifies the notion of ``gnosis" against the prejudice against it. As Christ points out, the ``legalists" are those who take away the key of gnosis. On the other hand, there is a gnosis that ``puffs up", and another that is ``knowing as we should." (1 Corinthians 8:1,2)

The issue is simple: gnosis is for the few (1 Cor. 13:8-10), it is the ``meat" of doctrine, not the ``milk". Hence, it will be hidden from most, who then will plausibly deny its existence. St Thomas Aquinas defines gnosis this way:

\begin{quotex}
As God works miracles in corporeal things, so also he does supernatural wonders above the natural order, raising the minds of some living in the flesh beyond the use of sense, even up to the vision of his own essence. 

\end{quotex}
\paragraph{Clement}
Smith then refers to \textbf{Clement of Alexandria} as the ``expositor par excellence of gnosis", since he embodies both the wisdom of the Greeks and the gnosis bestowed by Christ. Clement refers to an \emph{oral} tradition that came to him directly through the Apostles (\emph{pace} Sola Scriptura). Hence, Clement admits that some things cannot be written down.

Clement describes the qualities of the ``perfect man", which bears comparison to the ``true man" of Taoism. Such a man is ``impassible", i.e., transcends the movements of pleasure and pain, fear and desire, and so on. He is in a state of ``gnostic love".

Clement describes a conversion, a \emph{metanoia}, that happens in the ``twinkling of an eye". The conversion is to a divine vision characterized by the negation of duality. Smith draws out the consequences of nonduality. Nonduality is not a ``monism" like the night in which all cows are black. Rather gnosis is to be united to God while ignorance or disbelief is to be separated and divided.

\paragraph{Vedanta}
Besides Christian gnosis, Smith acknowledges three other traditions of gnosis: \textbf{Kabbalah}, \textbf{Tantrism}, and \textbf{Vedanta}. We discussed Tantrism earlier when reviewing Oscar Hinze. The Vedanta, according to Smith, constitutes the oldest \emph{living} tradition of gnosis. Therefore, there are ultimately two choices for doctrinal gnosis: the Christian and the Vedantic. Doctrinal gnosis is what can be spoken about, bearing in mind that true gnosis is beyond words and thoughts.

Without dismissing the Vedanta, Smith denies that Christian gnosis is merely a form of the Vedanta. Rather, it likewise is nondualist, yet the Trinitarian aspect adds a new dimension that is not reducible to the metaphysics of the Vedanta. Nevertheless, the encounter with the Vedanta is providential. The goal of that encounter is not to borrow from Eastern sources, Rather it is an anamnesis, a remembering of what had been forgotten.

The Kabbalah embodies the oral transmission pertaining to the Mosaic Revelation, providing access to the esoteric sense of the Torah. Smith devotes two chapters to the Kabbalah, concluding with its Christian development.

\subsection{Cosmology and Christian Kabbalah}

\begin{quotex}
None of the things which are comprehended by the senses or contemplated by the mind really subsist; nothing except the transcendent essence and cause of all. \flright{\textsc{St. Gregory of Nyssa}}

As fire is concealed by smoke, as a mirror by dust, as an unborn babe by the womb, so is Knowledge concealed by ignorance. \flright{\textit{Bhagavad Gita} III.38}

\end{quotex}
So far, \textbf{Wolfgang Smith} focused on the beginnings and the justification for a Christian Gnosis. Later on, he will deal with \textbf{Jacob Boehme} and \textbf{Meister Eckhart}, culminating in the notion of a ``Trinitarian Non-Dualism" vis-à-vis the Vedantic conception.

The middle chapters of Smith's book prepares the ground for a deeper understanding of Creation, the Incarnation, and the Trinity. First of all, Smith explains his cosmology. We won't need to linger on this, since it has been thoroughly covered in more detail in his books on the metaphysics of science. His contribution to the field is his interpretation of quantum mechanics that is consistent with traditional cosmology.

Smith then puts himself in the stream of the Hermetic tradition by reliance on the Christian Kabbalah. He explains the insights of Pico della Mirandola\footnote{\url{http://plato.stanford.edu/entries/pico-della-mirandola/}}, Johann Reuchlin\footnote{\url{https://en.wikipedia.org/wiki/Johann_Reuchlin}}, and Egidio di Viterbo\footnote{\url{https://en.wikipedia.org/wiki/Giles_of_Viterbo}} (``who ranks among the greatest Christian Kabbalists of all time"). Along the way, he mentions Friedrich Christoph Oetinger\footnote{\url{https://en.wikipedia.org/wiki/Friedrich_Christoph_Oetinger}}, the follower of Boehme, Valentin Weigel and the ``great" Paracelsus\footnote{\url{https://en.wikipedia.org/wiki/Paracelsus}}.

Rather than repeating Smith, which would be little more than a review of his reviews in any case, we will explain Smith's vision in slightly different terms. Thus it will be more like a jazz piece. The chord changes are the same, the melody is recognizable, we follow the same drummer, but we hope it will be a smoother composition. For Gornahoor readers, we can dispense with the intellectual justifications, but instead recommend the original text for that.

\paragraph{Anthropic Realism}
Smith calls his philosophical position \textbf{Anthropic Realism}: ``realism" because both universals and particulars are mind independent, and ``anthropic" because it stands on the bedrock of unmediated apperception. He admits his admiration for the phenomenology of Edmund Husserl, although Smith goes far beyond Husserl.

Scientific realism leaves no place for man, for the transcendent Intellect, which is a witness to the cosmos. It is not traditional cosmology that is absurd; it only seems absurd because scientific cosmology is absurd. In other words, scientific realism leaves the knower out of the picture. Of course, this goes back to \textbf{Arthur Schopenhauer} who asked the same question. Or even to \textbf{Bishop George Berkley}: does the cosmos exist if there is no one there to experience it? Smith's answer is a resounding ``No" and I'm sure that if you think about it, you, too, will feel compelled to agree.

The cosmos exists for us as an object of intentionality, i.e., man \& cosmos, microcosm \& macrocosm are complementary polarities. If one goes, so does the other. Certain artificial conundrums of naïve realism and Cartesianism are dissolved by anthropic realism. Viz., we don't see the image of a mountain or a tree in our brain; we actually ``see" the mountain or tree. The Cartesian universe, devoid of qualities, cannot be known and hence does not subsist on its own.

By distinguishing the corporeal world of apperception from the physical world without qualities, the existence of higher planes makes sense again. Above the quotidian corporeal world of our sense experience, there is the subtle or astral world, and the celestial world above that. No drugs or techniques can take us from the astral to the celestial world. Most today mistake the subtle as the ``spiritual" life. Modern man is stunted, with the higher centers merely virtual. Kundalini energy remains coiled up in the lowest chakra. Of course, gnosis awakens the higher centers, ``realizing" them, or making them actual rather than virtual.

\paragraph{Theophany}
Smith then deals with the exoteric teaching of Creatio Ex Nihilo. This view sees creation as static, as ``other than God". A breakthrough is possible to arrive at a nondualist conception: a \emph{theophany} rather than a creation.

Created things exemplify ideas in the Mind of God, which is ultimately the Divine Essence. God is simple, thus they are ``one" in God. Creation therefore exemplifies in its multiplicity what in God is ``one", i.e., nondual. But this is not a ``one time" affair. Creation is not manifested archetypes, but rather manifesting archetypes; this is a change from a static to a dynamic conception. We ``see" the thing, but we ``know" the archetype. That is, the archetype is ``repeated" in our mind, but not as another copy. Rather, it involves the participation of the human mind in the Mind of God.

Bottom line: There is no static creation, in addition to, and apart from, God. Rather there is God revealing himself. We get a sensual representation of the divine ideas; the world becomes phenomenal. Since God does not have ``sense" organs, God is conscious of the world through the human consciousness.

\paragraph{Kabbalah}
Smith devotes two chapters to the Kabbalah. Since there is extensive literature on the topic, we will just present some highlights. At some future point, a comparison with \textbf{Valentin Tomberg} and \textbf{Vladimir Solovyov} would be useful. Smith claims that, in addition to Christian gnosis, there are three other such traditions: Kabbalah, Tantra, and Vedanta. Due to its Scriptural roots, Smith asserts that the Kabbalah — of the three — is the most compatible with Christian Gnosis.

Smith finds in the Kabbalah keys to the Trinity, the Incarnation, Wisdom, the Eternal Feminine, and the four worlds, among other topics. The origins of Kabbalah lie in the mystical visions of Ezekiel. Tomberg does likewise, interpreting Ezekiel's visions as a revelation of the archetypal world and the chariot as the central symbol of cosmic knowledge.

The En-Sof is the Absolute and Infinite Godhead and the ten Sefiroth are God's self-revelation. It is the world, yet the highest Sefiroth represent the Trinity:

\begin{itemize}
\item \textbf{Keter}: the transcendent 
\item \textbf{Hokhmah}: Wisdom 
\item \textbf{Binah}: Intelligence 
\end{itemize}
There is then an emanation down to Malkuth, or the World. Malkuth is also represented by the Feminine principle of Shekinah, which in turn is tied to Binah, the Holy Spirit.

The mystery of sex is symbol of the love between the divine I and the divine You. The \emph{hieros gamos} is the central fact in the whole chain of divine manifestations

Adam and Eve, man and cosmos originate in higher world. Interestingly, Smith comes close to the ancient Gnostic idea of the world, as we now experience it, as the creation of a ``Demiurge", an inferior being. Originally, the material world which we behold with our senses and contemplate with our rational mind did not exist. It was ``sin", the awareness of evil, that created this world. This is not precisely a ``new world" but more like a new way of knowing. Adam becomes attached to the glamour of the world, worshiping the Shekinah exclusively, and failing to see the connection to the higher Sefiroth.

There is a deep truth here, since it is a commonplace that God ``created the world". But are mosquitoes, cystic fibrosis, and so on really part of the original creation? Clearly the answer is No, otherwise the ``Fall of Man" and the ``expulsion from the garden" would make no sense.

This bears connection to Tomberg. Adam's originary orientation was vertical, that is, he experienced the theophany directly. The fall, then, was to change orientation to the horizontal, the false belief that things are independent of any transcendence. So the fallen world was not created by another being, or demiurge, but rather by the Antichrist, which is an egregore. It is only ignorance that considers the egregore to be independently real.

The Kabbalah describes four worlds: emanation, creation, formation, and action. In Neoplatonism, these correspond to the One, the Intelligible, Soul and Matter. Leaving out emanation which is totally transcendent, these correspond to the angelic, subtle, and proto-material worlds. They are also the three worlds of the Hindu Triloka.

Tikkun, or redemption, is the restoration of the ideal order. Today, of course, that primal order is seldom understood. Its foundations are polarity and hierarchy. There is a hierarchical order to everything, from the spiritual, to the subtle, to the material. Manifestation requires a series of syzygies; this notion will be expanded in a further explanation of Tomberg's ideas on Salvation and Evolution.

The divine order is Hierarchy, which revolutionaries oppose. Moreover, without polarity, there is no manifestation. Hence, the attack starts with the flesh and promotes the idea that polarity is nothing more than an arbitrary preference or orientation.

\paragraph{Jesus Christ}
On the other hand, the Shekinah, as the Theotokos, can lead us upward out of ignorance, whether through the connection with Binah, or with the Tifereth, the Heart Sefiroth, which represents Jesus Christ. Smith points out that Christ has the perfect balance of masculine and feminine virtues. In a series of paradoxes, reminiscent of the Tao Te Ching, he explains:

\begin{itemize}
\item He is the King who rules, yet is meek and lowly in heart 
\item He commands, yet is submissive to all 
\item He the Judge of kings and the Friend of sinners 
\end{itemize}
In short, he personifies the Word [Logos] that was in the beginning. In the Chinese translation of the Gospel of John, ``Tao" is used to translate ``Logos".

\subsection{Meister Eckhart}

\textbf{Carl Jung} regarded \textbf{Meister Eckhart} and \textbf{Dante} as two of the Ten Pillars of the Bridge of the Spirit\footnote{\url{http://jungiancenter.org/wp/ten-pillars-bridge-spirit/}}.

\begin{quotex}
\emph{Faust} is the most recent pillar in that bridge of the spirit which spans the morass of world history, beginning with the Gilgamesh epic, the \emph{I Ching}, the \emph{Upanishads}, the \emph{Tao-Te-Ching}, the fragments of \textbf{Heraclitus}, and continuing in the \emph{Gospel of St. John}, the letters of \textbf{St. Paul}, in \textbf{Meister Eckhart} and in \textbf{Dante}… \flright{\textsc{Carl Jung}}

\end{quotex}
The Traditional author, \textbf{Ananda Coomaraswamy} (AKC), in \textit{Vedanta and the Western Tradition}\footnote{\url{https://gornahoor.net/?p=556}} exclaimed:

\begin{quotex}
\textbf{Eckhart}, with the possible exception of \textbf{Dante}, can be regarded from an Indian point of view as the greatest of all Europeans. 

\end{quotex}
\textbf{Valentin Tomberg} included Meister Eckhart among those who arrived at the point of contemplation using Thomism as their starting point. In the Letter on the Sun, we read:

\begin{quotex}
Meister Eckhart, Ruysbroeck, or, lastly, St. John of the Cross are spirits amongst whom you will search in vain for a spirit of opposition to scholasticism. For them also it was true that scholasticism was ``like straw", but they knew from their own experience that this straw proved to be an excellent combustible. They certainly surpassed scholasticism, but they did so by attaining its aim. For the aim of scholastic endeavor is contemplation, and it is mysticism which is the fruit of the scholastic tree. 

\end{quotex}
The starting point of scholasticism are the two sources of knowledge: Reason and Revelation. While Reason strives to understand the nature of the world, it can only penetrate part of the way to knowledge of the Divine Essence. For the latter knowledge requires Revelation. However, deep contemplation seeks to penetrate into an intuitive grasp of the Divine. According to Tomberg, this task is also the Hermetic task, viz., to unite intelligence and the intuition of faith.

Among those who are ``very advanced in substantially" uniting them, there is \textbf{Nicholas Berdyaev}. Berdyaev, in \emph{The Divine \& the Human}, asserts that the opposition of the divine and the human will be overcome in the Holy Spirit. Without going into all the details at this time, he claims that the era of the Spirit was anticipated by Meister Eckhart, then by his successors \textbf{Johannes Tauler}, \textbf{Jacob Boehme}, and \textbf{Angelus Silesius}.

\textbf{Wolfgang Smith}, in \emph{Christian Gnosis}, takes up the same themes: the union of the divine and the human, as well as the role of the Holy Spirit. Through an analysis of Eckhart's mystical vision, Smith describes a Trinitarian non-dualism. There are two phases: the Trinity and the birth of the Logos in the soul.

\paragraph{Trinity}
Eckhart develops ideas from \textbf{Dionysius} and \textbf{Augustine} in his understanding of the Trinity. Augustine regarded the Father, Son, and Holy Spirit as analogous to ``Being, knowing, and loving or willing." \textbf{Vladimir Solovyov} thought highly of Augustine's explanation of the Trinity from the \emph{Confessions}. Augustine related the Trinity of Plotinus to the Christian understanding as shown in this chart:

\begin{table}[h]\centering\scriptsize
\begin{tabular}{ccc}\toprule
\textbf{Plotinus} &
\textbf{Augustine} &
\textbf{Trinity}\\
The One &
Being &
Father\\
Idea-world &
Knowing &
Son\\
Psyche &
Willing (or loving) &
Holy Spirit\\\bottomrule
\end{tabular}
\end{table}
Augustine derived this understanding from the phenomenology of his own spiritual life. Solovyov shows how these three acts are identical to each other.

\begin{enumerate}
\item I am the One who Knows and Wills 
\item I know (or am conscious of) my Being as well as my Willing. I know that I am (being) and that I will. 
\item I will myself as one who is and who knows. 
\end{enumerate}
Each of these acts of the Spirit is completed by the other two, so that they constitute a unity.

The Father is Being, even beyond being, the One, the unknowable essence of the Godhead, or Thearchy as Dionysius called it. The Father is beyond sense experience and intellect so that the self-emptying of apophatic theology is ultimately the only path to approach God.

He knows himself in the Son, the articulated word. The Son, as the Logos, conveys the intelligibility of the Father. The Father knows himself as reflected in the Son, and the Son knows the Father. The Holy Spirit is the unifying love.

This is the way to know God. Eckhart makes this clear:

\begin{quotex}
if you would know God, you must not merely be like the Son, you must be the Son yourself. 

\end{quotex}
It is not enough to be ``Christ-like"; rather we must be transformed. When we put on the mind of Christ, we become the Son, so that God is reflected in the soul. The Holy Spirit unites the divine and the human.

\paragraph{The Birth of the Word in the Soul}
There is difficulty for the hylic personality to attain to spiritual vision, since sensual imagery is his primary mode of knowing. Eckhart writes:

\begin{quotex}
Some people want to look upon God with their eyes, as they look upon a cow, and want to love God as they love a cow. 

\end{quotex}
So Jesus is the Son incarnated in the flesh. Nevertheless, for the Holy Spirit to come, he must go:

\begin{quotex}
It is expedient for you that I go away: for if I go not away, the Comforter will not come unto you; but if I depart, I will send Him unto you. (John 16:7) 

\end{quotex}
The point is clear. In Jesus Christ, the human and divine were united in a visible way. For us, however, the Spirit must come in order to overcome that division that separates us from the divine in our own consciousness. Eckhart warns us about the neglect of the Spirit. Although Christ ascended into Heaven and sits at the right hand of the father, that avails us nothing.

\begin{quotex}
What would it avail me if I had a brother who was a rich man, and for my part I were a poor man? What would it avail me if I had a brother who was a wise man, and I were a fool? 

\end{quotex}
He continues:

\begin{quotex}
The Heavenly Father brings forth his only-begotten Son in Himself and in me. 

\end{quotex}
The birth of the Logos in the soul\footnote{\url{https://www.meditationsonthetarot.com/the-word-is-made-flesh}} is the whole point of \textbf{Letter II: The High Priestess} in \emph{Meditations on the Tarot}. The Son is by generation, but the Son who lives in us, makes us children of God by adoption. This is what the Meditations calls Initiation. Eckhart goes even further, bringing out the consequences on his non-dual understanding. In Sermon 6, we read:

\begin{quotex}
The Father begets his Son like himself in eternity. ``The Word was with God and God was the Word." It was the same in the same nature. I will say more: he has begotten him in my soul. Not only is it with him and he with it alike, but he is in it. The Father begets his Son in the soul in the same way as he begets him in eternity, and not otherwise. 

\end{quotex}
The sacrifice of the Eucharist is not a repetition nor a re-enactment, but is the same sacrifice as Calvary. This is something that happens in eternity, not in time. Similarly – and this is what Eckhart is driving at – the birth of the Logos in the soul is not a different ``logos" but is the same Logos that is with the Father eternally. That is how it looks from the standpoint of eternity.

It is the birth of the Logos in the soul that brings knowledge of the Father, otherwise it would be impossible. That is the non-dualism. This second birth overcomes the duality of wills. Sermon 12 states:

\begin{quotex}
A person who is so established in the will of God wants nothing else but what is God and what is God's will. 

\end{quotex}
The Father knows himself in the Son, and the Son knows the Father. When ``I" am united with the Son, we can understand Eckhart:

\begin{quotex}
The eye in which I see God is the same eye in which God sees me. My eye and God's eye are one eye and one seeing, one knowing, and one loving. 

\end{quotex}
And in Sermon 6:

\begin{quotex}
In the innermost spring I went forth in the Holy Spirit. There is one life and one being and one activity there. 

\end{quotex}

\subsection{Jacob Boehme}

\begin{quotex}
I was a hidden treasure, and I desired to be known; therefore I created the creation in order that I might be known. \flright{\textsc{Hadith Qudsi}}

Origen, Dionysius the Areopagite, \textbf{Jacob Boehme}, Louis Claude de Saint-Martin, Vladimir Solovyov and Nicolas Berdyaev, for example, show in their works a progress which is very advanced in \emph{substantially }bringing together intelligence and the intuition of faith. \flright{\textsc{Valentin Tomberg}}

Jacob Boehme has to be termed the greatest of Christian gnostics. The word gnosis I employ here not in the sense of the heresies of the first centuries of Christianity, but in the sense of knowledge basic to revelation and dealing not with concepts, but with symbols and myths; contemplative knowledge, and not discursive knowledge. This is also a religious philosophy or theosophy.
\flright{\textsc{Nikolai Berdyaev}, \emph{Studies Concerning Jacob Boehme}\footnote{\url{http://www.berdyaev.com/berdiaev/berd_lib/1930_349.html}}}

\end{quotex}
In \emph{Christian Gnosis}, \textbf{Wolfgang Smith} reveals himself as a follower of \textbf{Franz von Baader} in the stream of Christian theosophy. That is why he concludes his work with the ideas of \textbf{Paracelsus}, \textbf{Jacob Boehme}, and \textbf{Meister Eckhart}. Smith interprets the latter figure as the sources of a distinctively Christian Trinitarian nondualism, not inferior to the nondualism of the Vedanta. This post will deal with Smith's chapter on Jacob Boehme, and the final one with Meister Eckhart's Trinitarian nondualism.

\textbf{Jacob Boehme}\footnote{The main source used for Boehme's teachings is \textit{Five Christian Principals} by Rene Cossey.} is the model of the esoterist for the coming Age: He was a layman — a cobbler by trade, just like my grandfather — was married and raised a family. Yet behind that quite ordinary looking exterior, he was undergoing a profound spiritual transformation. He tried to live an obscure life; his first book was distributed just to his friends. Toward the end of his life, he was finally able to dedicate his time solely to his spiritual writings, thanks to some wealthy patrons.

\paragraph{Excursus on Gnosis}
Berdyaev describes the foundation of Boehme's gnosis in the following paragraph:

\begin{quotex}
We ought to be re-united with the traditions of the theosophy and anthroposophy of J. Boehme, in truth with a Christian theosophy and anthroposophy. And moreover, even more deeply ought we to be re-united with the traditions of the esoteric, hidden Christianity. But the fruition of the great traditions of Boehme and of Christian gnosticism ought to be creative, it ought to guide along the path of a completely new, creatively-active knowledge. Modern people, seeking God and the Divine life, are very afraid of thought and knowledge, and the basic thrust of their will often becomes anti-gnostic. They admit the possibility only of a passive, abstract knowledge. They cannot accept knowledge as a creative act, bearing life into the light of the world, of knowledge as being and life.  \flright{\textsc{Nikolai Berdyaev, }\emph{Theosophy and Anthroposophy in Russia}\footnote{\url{https://www.berdyaev.com/berdiaev/berd_lib/1916_252b.html}}}

\end{quotex}
The fundamental notion is that gnosis is creative. Discursive knowledge, on the other hand, is passive and abstract, i.e., it is something to be absorbed as is. Gnosis is the knowledge of the heart. It is creative in the sense that the knower needs to raise his level of being, and that is a creative process. In other words, he needs to be the artist of his own life, his own soul; only in that way can he hope to achieve gnosis.

The next notion is that Boehme has both a cosmology and an anthropology. They are intimately related as the macrocosm to the microcosm. By knowing himself, man knows the macrocosm, and by knowing the macrocosm he knows himself. This was more fully developed by Boehme's follower Johann Gichtel\footnote{\url{https://gornahoor.net/?p=8618}}, who showed that the planets are related to the inner states or chakras.

Of course, this assumes a traditional cosmology that involves the following degrees of existence:

\begin{itemize}
\item Mineral or physical 
\item Plants 
\item Animals 
\item Intellect 
\item Celestial or Angelic 
\item Divine 
\end{itemize}
Those states of being also exist within man, even if most people live their lives at the levels of plant and animal life, with excursions into the truly human state or perhaps higher states. As a matter of fact, it is considered a marker of high intelligence and good education to believe that only the mineral or physical state of being actually exists. Thus, all interior life is believed to be the result of electro-chemical processes. For them, I have selected this theme song\footnote{\url{https://www.youtube.com/watch?v=dkSfdGNuD0c}}.

The following sections briefly summarize Boehme's system in order to understand Smith's interest. As we pointed out, this should not be read as passive, abstract and discursive knowledge to be passively received. Instead, make the effort to re-experience Boehme's vision in your own consciousness.

\paragraph{The Ungrund or Abyss}
Boehme has his own particular interpretation of Traditional metaphysics. In common with them, he starts with the unmanifest source or all, describes the Law of Three that makes manifestation possible, and then the Law of Seven indicating the stages of that manifestation.

The Ungrund desires to reveal himself through manifestation. The universe is the outcome and development of One Grand Thought. All things are governed by one central law and all planes of existence are related. This is the Law of Correspondences.

The Abyss contains within itself everything and nothing, i.e., everything potentially but nothing manifestly. Within the Abyss is an eternal uncreated Will (Byss). It desires to become manifest, to be something.

The Will fashions a Mirror which reflects all things, making them manifest. The Mirror is Eternal Wisdom, Eternal Idea, Virgin Sophia, and the Infinite Mother, while the Will is the Infinite Father. Like the Tao, Boehme describes this relationship:

Abyss $\Rightarrow $ Duality $\Rightarrow $ Trinity.

The Father-Mother begets a Son; these energies are diffused by the Holy Spirit. Through the union of the Will and Wisdom, the unmanifest becomes manifest, the latent becomes active.

In Guenon's system, the Ungrund is Nonbeing, and Being is its first principle. Boehme claims, moreover, that the Ungrund is God as he is in himself. The Ungrund makes the Ground, i.e., Being, which is the Father. This is similar to the Palamite understanding assuming that the Ungrund is the unknowable essence of God, and the Father is the energies, or the Act of Being as in Thomist metaphysics. The Logos concentrates those energies, which are diffused through the Holy Spirit. Thus, the generation of all things takes place and the unmanifest becomes manifest.

Obviously, this should not be understood temporally, since it is all prior to time itself.

\paragraph{The Ternary}
In \emph{Symbolism of the Cross}, \textbf{Rene Guenon} describes the three gunas, which are ``essential, constitutive and primordial qualities or attributes of beings envisaged in their different states of manifestation". Boehme likewise recognizes those qualities. There are two which he calls Fire and Light. These contrasting principles exist in all things. The third manifests in our external nature. This can be summarized as:

\begin{itemize}
\item Fire, Wrath, Law, Dark Principle, latent and unmanifest (nonbeing) 
\item Light, Mercy, Love 
\item Manifestation or Being 
\end{itemize}
The Dark Principle is unmanifest, since it would be understood as the wrathful or vengeful aspect of God. Nevertheless, Man has free will, so he can will the dark principle. A contrast is necessary for conscious awareness. This is not a problem provided it is transmuted by the Light or Love Principle. Otherwise, man is in the ``false imagination", resulting in a misunderstanding of God and the World.

In other Traditions, these principles are called

\begin{itemize}
\item Brahma the Creator/Shiva the Destroyer, 
\item Ormuzd/Ahriman, 
\item God/Devil 
\end{itemize}
Vishnu, then, is the Preserver, balancing the two principles. This requires knowledge and wisdom gained through experience. This reconciles Law and Love. By following Truth, man transmutes evil into good, otherwise he becomes a slave to evil.

The Polarity occurs on every plane of manifestation. It is reflected as

\begin{itemize}
\item Positive/negative 
\item Active/passive 
\item Masculine/feminine 
\item Action/reaction 
\end{itemize}
This Polarity — or Sex — is the law of all manifestation, the creative power of the universe. That is the cosmic order and any deviation from it is disordered.

\paragraph{The Septenary}
Boehme says there are seven properties or forces through which the Divine energy operates. Each quality has its own essence, forming together one harmonious whole. They permeate all manifest things. Although Boehme talks about them as if sequential, they operate simultaneously, outside of time. These are the seven qualities, followed by their names:

\begin{itemize}
\item \textbf{Contraction}: The desire drawing all towards itself (\emph{Desire}) 
\item \textbf{Friction}: expansive force creating a dual action and differentiation. Desire going into multiplicity (\emph{Motion}) 
\item \textbf{Sensibility}: rotary motion of 1 and 2. The wheel of life. (\emph{Sensation}) 
\item \textbf{Lightning or Fire}: The Spirit diffuses a mild light, transforms the dark principle, ending strife 
\item \textbf{Love or Light}: love binds, builds up and harmonizes the principles of joy 
\item \textbf{Audibility}: intelligible or vital sound. The manifestation of Life. (\emph{Intelligence}) 
\item \textbf{Essential Wisdom}: gathers the previous six into one harmonious whole 
\end{itemize}
Contraction, Friction, and Sensibility form the ``Dark Principle". Contraction is the desire drawing things to itself; it is harsh, cold, and sharp. It is a kind of magnetic attraction, congealing Nothing into Something. Friction is expansive, thus causing differentiation. It is desire becoming multiplicity. Sensibility is the wheel of life brought about by the action and reaction of Contraction and Friction. This wheel is experienced as wandering and anguish, which amasses itself into an Essence. In other words, the first principle is Attraction, then Repulsion, resulting in Circulation.

The Lightning Flash is the eruption of the Spirit, who transforms the Dark principle by the suffusion of Light. This ends the Strife of Attraction and Repulsion. This is the beginning of Consciousness and Life.

Love then builds up and harmonizes the principles, bring joy and perfection. Audibility brings intelligence. Finally, Wisdom brings harmony and bliss.

This process is the ``Great Work", the \emph{Magnum Opus}, that must be performed in each of us in order to reclaim our original nature and ``the great transmutation which Christ Jesus accomplished in His Own Person."

\paragraph{The Created Universe}
The Uncreated Heaven is perfect and complete and is part of God. However, the Created Universe is outside and apart from God. Therefore, it needs to develop from incompleteness to completeness. This does not exhaust God, who is Infinite, but rather enriches him.

God produces from his own eternal nature and wisdom, in which all things exist as possibilities of manifestation. Created things manifest successively in time. Creation is in Equilibrium, the harmony of the active forces (Light) and passive forces (Dark).

Creation unfolds through endless circles. Outside the circle of the Uncreated heaven, there is the created heaven or angelic world. There are three hierarchies headed by Michael, Lucifer, and Uriel (corresponding to the Father, the Son, and the Holy Spirit).

\paragraph{Fall of Lucifer}
Although the Dark Principle is passive and unmanifest in God, and never becomes active, the possibility of evil arises in free beings. Free beings can choose to be centered in the natural egoic Self. This is to remain at the level of Contraction. Or else they can be centered in the Light Center, the expansive power of Love. Evil and disorder enter the world when creatures become centered in that lower self. Once selfish desire is kindled, it just becomes stronger. Instead of letting the Divine Light reconcile opposing forces, they get stimulated resulting in anguish. Life becomes struggle and anxiety.

Evil first entered the Cosmos when Lucifer, the head of our universe, became self-centered thereby rejecting the divine order. Lucifer's fall was the result of pride and the pursuit of knowledge apart from God. He rejected the feeling of humility in the face of the Divine mysteries. Curiosity and the desire for novelty became dominant. Thus evil, Hell, and the Dark Principle entered the world.

Our personal task, therefore, is to reconcile the extremes of the forces of Darkness and Light. As the Hermetists taught, that is necessary for the recovery of the Prima Materia, the element of immortality. Then it can be molded by the Will for creative purposes, without painful toil, ``transforming work into play".

This must begin with self-knowledge in order to recognize the Luciferian temptations in our own souls. This should be easy to do, since Lucifer is the Prince of this World.

\paragraph{Adam's Fall}
Earth was situated spiritually within the sphere of Lucifer. Lucifer split the world into Darkness and Light, separating God's Wrath from God's Love, making Equilibrium difficult. The world became dense and gross.

As we pointed out above, man is not just a ``rational animal", but rather a composite being containing elements of all things. Thus he has facilities to know things on all planes, including the physical, astral, physical, and divine things.

The human race was ethereal, so Adam was a luminous being permeated by a celestial essence. His mind was innocent, no knowledge of evil, without avarice, pride, envy, or anger. Adam's body was not dense, since the inner life was the master of the body. Adam's inner life was in touch with Heaven; he could communicate with God, the angels, and nature. His being was tripartite, in harmony with each other:

\begin{itemize}
\item \textbf{Spirit}: light principle 
\item \textbf{Soul}: dark principle 
\item \textbf{Body}: in the world of sense formed through the union of the light and dark principles 
\end{itemize}
Nevertheless, virtue requires the possibility of vice and the experience of temptation. So Adam was tempted by Lucifer, decentering his attention onto the world. He was tempted to experience the pairs of opposites, thus disturbing his previous state of Equilibrium. This was the Fall into duality: the experience of each side of light and dark in isolation.

Hence, our task to reconcile these opposing principles. Unlike Lucifer's Fall, Adam did not really wish to oppose God, but merely to experience earthly pleasures. Therein lies the clue to the way back.

\paragraph{Adam and Eve}
Originally, Adam was a dual unity, including Eve. Boehme describes Adam-Eve like this:

\begin{quotex}
Adam was a man and also a woman, and yet neither of them distinct, but a virgin full of chastity, modesty, and purity, namely the image of God.

\end{quotex}
When Adam's attention was enticed by the world, Eve (the feminine part of his nature) prompted him to separate. Noting that the animals were male and female, he/they were overcome with desire to copulate like mammals\footnote{\url{https://www.youtube.com/watch?v=xat1GVnl8-k}}.

Adam's deep sleep refers to his forgetfulness of the angelic world; he then awoke in the external world, separated from Eve. He lost the unitive consciousness in dualistic thinking. They ate the fruit of the Tree of the Knowledge of Good and Evil, becoming subject to death, decay, and corruption. Now part of the sidereal universe (i.e., under the stars), they lost awareness of anything higher.

Because of Lucifer and Adam, even nature has become gross. In the world there is a mixture or alternation of the two principles of the Dark and the Light. There are storms and beautiful weather; poisonous plants and delicious fruit; savage beasts and noble animals. Curse, decay, corruption, and death struggle with blessing, health, and life.

Man has sunk to the animal plane of existence, tending towards the bestial. His inner disposition is like the animals. Man must extirpate these instincts to become wholly human, and restore the image of God. Adam should have restored creation to order, following the fall of Lucifer. Unfortunately, he is unable to bring the opposing forces into Equilibrium. Christ, the second Adam, re-establishes man in his primal dignity as Lord of Creation.

\paragraph{Atonement and Redemption}
Jesus was a Divine Human Being with the essence of both God and Mary; he combined the highest part of the masculine (forcefulness) with the best attributes of the feminine character (tenderness). Christ is the Son of Virgin Mary, but also of the Heavenly Virgin who united herself with Mary.

Whereas Adam fixed his imagination on the lower world, Jesus fixed his imagination wholly upon the Father subjecting the lower principle to the higher, re-establishing equilibrium. The resurrection body is paradisiacal body, which eventually disappeared (ascended to Heaven). Heaven is not a distinct place, but interpenetrates the physical universe. This is hidden from normal sight.

Regeneration by the Holy Spirit is required for Equilibrium. This requires and inward transformation, not just an historical faith in Jesus. Man must forsake the principles of Darkness and the world (lust and appetites). The Light principle was lost through Adam but regained with Christ. Prayer is the means to soar above the center of nature. Will and desire become one. It is necessary to be in harmony with the Divine Will, and then conscious union and knowledge of God.

On earth, the soul is already either in heaven or hell. So it does not ``go" anywhere after death, yet the soul is fixed at the moment of death and cannot be changed. However, souls in a half-regenerate condition can ultimately reach heaven after a period of purgation.

Creation was the act of the Father, the Incarnation was the act of the Son, Holy Spirit will bring about the end of the world. This will bring about nondual awareness. This is reminiscent of Joachim of Fiore.

\paragraph{Nondualism}
Wolfgang Smith, then, interprets Boehme's theosophy in a nondual way. Christ is the Incarnate God who reveals himself. This happens on several levels:

\begin{itemize}
\item On the plane of eternal nature as the summation of the septenary cycle 
\item On the plane of human history as Jesus of Nazareth 
\item God becomes incarnate in all who are born in Christ 
\end{itemize}
The first is creative and the second is redemptive, i.e., the restoration of a state which had been lost. Man and cosmos are in flux, a ``kind of samsara", which will have its fulfilment in the Kingdom of God.

Smith concludes this chapter with three questions:

\begin{itemize}
\item Does Boehme's doctrine differ from the Kabbalah? 
\item Is it Trinitarian? 
\item Is it Nondualist? 
\end{itemize}
The answer to the first question is a qualified Yes, with some adjustments. But what is more interesting it the idea of a Trinitarian Nondualism, which Smith will develop more fully in the final two chapters on Meister Eckhart. Specifically, Deliverance is not a matter of just knowing the transcendent Father, but of knowing both the Father and the Christ. This is one act of knowing, not two, and is incarnational as well as transcendent. The Mystical Union is the Supreme Identity or, as Boehme describes it:

\begin{quotex}
What is void of will is one with the Ungrund … which is God himself.

\end{quotex}
\paragraph{Postscript on Personal Development}
I've avoided discussing what is most discussed: e.g., the life and sanctity of Boehme, his siddhis, the sources of his doctrines, or whether they are heterodox. Ultimately, they need to be comprehended, if at all, internally, in one's conscious life.

This is clearly difficult because of the falls of Lucifer and Adam. We have lost our birthright gift of understanding transcendence due to our attachment to the glamor of the world process. These are some exercises in conscious awareness that you can try at home.

\begin{itemize}
\item Lucifer exchanged an intellect centered on God to one centered on the world. Cunning, guile, idle curiosity, and the like, replaced genuine creativity. See \textit{The Intellectual Center}\footnote{\url{https://www.meditationsonthetarot.com/intellectual-center}} for the distinction between the higher and lower functions of the intellect. 
\item Adam rejected a vision of wholeness for dualist knowledge and attachment to the natural world. He entered a world of Strife instead of one of Love, all the while assuming he can somehow remain on the ``pleasant" side of duality. 
\item In that state, it is assumed that only the physical world exists and all problems have a scientific or material solution. The very idea of a quest for higher consciousness is rejected. 
\item By his attachment to sexual reproduction and its ersatz simulations, he finds it difficult to refocus attention upwards. The same sexual energy could be used to achieve the original wholeness, or alchemical marriage, or even like Dante and Beatrice, to achieve the mystical union. 
\end{itemize}
These aspects of consciousness need to be experienced personally. Unless there is an awareness of being in a fallen state, then the very idea of redemption makes no sense.

\flrightit{Posted on 2016-07-20 by Cologero }

\begin{center}* * *\end{center}

\begin{footnotesize}\begin{sffamily}



\texttt{Boreas on 2016-07-20 at 22:19 said: }

This is a great book dear Friend. It's great you've recently taken Wolfgang Smith's opuses in deeper examination. Was it me who first recommend it to you? I have given this book two times from my bookshelf to dear friends, as a Christmas present and so on. I've read and bought all Susijengi Sepän (Wolfgang Smith's) books and they have been life's changing lately. I loaned three of his books to psychologist-fyscicist Friend of mine and he became interested. Well, to cut the story short, namaste.


\hfill

\texttt{Cassiodorus on 2016-07-20 at 23:50 said: }

Colegero,

I an hoping that you will be discussing Smith's Trinitarian nondualism. I am not sure- but I think it is not actually in alignment with Guenon and Schuon's metaphysical point of view. It seems to me that Wolfgang Smith's position resembles that of Jean Borella's or perhaps Robert Bolton?


\hfill

\texttt{Cologero on 2016-07-21 at 12:50 said: }

@Cassiodorus, we will get to Smith's Trinitarian nondualism, although he attributes it to Eckhart, not himself. I don't know Borella's view on it, although Smith does mention him several times.

You are correct, it is a deviation from Guenon, yet not in the way of a refutation.

Keep in mind the dogmas like the Trinity used to be restricted to initiates, and were not intended for the general public.


\hfill

\texttt{Matt on 2016-08-12 at 10:54 said: }

Boehme, who will be the subject of the next part, has an interesting take on the Fall. There were two falls: the Fall of Man and the Fall of Lucifer. Lucifer's Fall lead to the natural evil/privation in the world, and Adam's Fall lead to the moral evil/privation in the world.

One doesn't usually find an examination of the extent to which Lucifer's fall affected the world beyond merely bringing temptations to the human being. I guess Boehme stands out in this regard.


\hfill

\texttt{Babbus on 2016-08-12 at 19:54 said: }

Mr. Salvo, here's a link to Corbin's \textit{The Paradox of Monotheism}: 

\url{https://www.scribd.com/doc/300003894/The-Paradoxe-of-Monotheism-Henry-Corbin}

I hope to have been of assistance.


\hfill

\texttt{Boreas on 2016-09-10 at 16:03 said: }

Excellent synopsis again, Cologero! I haven't been here for a while, but as a synchro I just yesterday evening read the chapter on the Gnosis of Boehme.

Boehme's explanation of the Fall of Lucifer is interesting. I have lately been thinking that it was a cosmological and ontological necessity to bring the fire of mind to man, which lead to the involutionary cycle with all its effects of cosmic evil, to the outgoing, which is followed by the evolutionary cycle, or the in-going and the path of return that was perfected in Jesus Christ, who ``redeemed Lucifer" and established the law of mercy and love as the highest standard in earth instead of the old covenant which was based upon justice. Any thoughts on this?


\hfill

\texttt{Barnabas Fiacabrili on 2016-10-07 at 02:06 said: }

Astounding. I personally want to thank you for showing me the truth of Christinaity. Years ago I left the church in favor of vendata. I know now, thanks to this website and all the great intellectuals I have read since, how inseparable they are. By grace, direct experience and knowledge will solidify this.

Thank you and God bless!


\hfill

\texttt{Victimarius on 2016-10-07 at 07:49 said: }

Excellent article. I would also like to thank you for it. Illuminating.


\hfill

\texttt{Mark Citadel on 2016-10-14 at 05:03 said: }

Astounding stuff, definitely has enhanced by knowledge of gnosis. Just unfortunate so many people associate it with the early century heresy. The section on the Fall of Lucifer in particular cleared up a few issues I was having with the Ahriman comparisons.

I really must get around to picking up your book!


\hfill

\texttt{Matt on 2016-11-16 at 18:32 said: }

Wolfgang Smith had a correspondence with Fr. Malachi Martin, a well-known Catholic priest in the West. Their correspondence covered topics on Tradition and its relation with the Church. These written correspondences were just recently released by Smith. A key focus was on Jacob Boehme. Martin appears to share Smith's high esteem of Boehme, as well as appreciating Hermeticism (Martin describes Hermetic teaching as pre-Christ Christology).

Below is a link to the first bit of their written correspondence where they discuss Boehme's teaching on Lucifer's fall. What I thought was most noteworthy is Smith and Martin's agreement that Boehme's insight on this topic renders an interesting interpretation of the creation account in Genesis. Smith has drawn out the implications of this in a number of his wiritngs and lectures, which have been touched upon in the previous Gornahoor posts in this series. Gornahoor readers and Cologero may find it interesting.

\url{http://www.christendom-awake.org/pages/book-promotions/in-quest-of-catholicity/book-extracts.pdf}


\end{sffamily}\end{footnotesize}
